Context-adaptive access control (CAAC) lets IoT and cyber-physical deployments trade
authentication strength, inspection depth and session freshness against energy and
mission utility. When the context channel is spoofed, the same mechanism becomes an
attack surface: a coherent ``benign'' context can talk the policy below the posture
that the physical process requires. This paper separates the part of adaptation that
may be optimised from the part that may not. We formulate CAAC as a constrained
partially observable control problem over five latent threat states, with a forward
Bayesian belief, cross-modal consistency weights, and an admissible set defined by
unconditional invariants for actuation and safety-critical traffic. The controller
optimises expected utility only inside that set; if the set is empty it executes a
precomputed per-operation fail-safe action, and the enforcement point independently
re-checks the floor. A 24-node, 2{,}000-epoch, 10-seed simulation across seven attack
scenarios holds invariant violations at 0 per 1{,}000 decisions for every
invariant-enforcing policy, versus 393 for a static-low policy and 327 for
unconstrained risk adaptation. The full controller attains 0.120 mean
attack-scenario utility against 0.084 for static-high and 0.077 for threshold CAAC,
while denying 0.029 of safety-critical requests against 0.088. Residual exposure
remains and is reported: the full controller grants 0.177 of malicious requests
against 0.127 for static-high, concentrated in an explicitly out-of-model scenario
in which the node's trusted execution environment is subverted and the belief
collapses (Brier 0.437, poisoning AUROC 0.273). Within the stated threat model the
attested device anchor and cross-node dispersion survive poisoning, and the
guarantee is carried by the unconditional invariants rather than by detection. A
trace-derived calibration against IoT-23 sets the benign-to-malicious mean-log-rate
contrast to 1.69.
